\documentclass[11pt, oneside]{article}   	% use "amsart" instead of "article" for AMSLaTeX format
\usepackage{geometry}                		% See geometry.pdf to learn the layout options. There are lots.
\geometry{letterpaper}                   		% ... or a4paper or a5paper or ... 

\usepackage[parfill]{parskip}    		% Activate to begin paragraphs with an empty line rather than an indent
\usepackage{graphicx}				% Use pdf, png, jpg, or eps§ with pdflatex; use eps in DVI mode
								% TeX will automatically convert eps --> pdf in pdflatex		
\usepackage{amssymb}
\usepackage{cite}

% numbered examples
\usepackage{gb4e}
\usepackage{enumitem}
\usepackage{cancel}
\usepackage{amsmath}

\usepackage{mhchem}
\usepackage{lewis}
%\usepackage{urlbst}
\usepackage{url}
\usepackage{hyperref}


% Scientific Laws
\newtheorem{definition}{Definition}
\newtheorem{law}{Law}
\newtheorem{theory}{Theory}
\newtheorem{hint}{Hint}

\title{Module 12 (Reprise)\\ Fun with Gases and a .44}
\author{Mark Peever\\ \texttt{mpeever@gmail.com}}
\date{March 6, 2026}

\begin{document}
\maketitle

\section{Overview}
This is a short side-quest.
Let's consider our Ideal Gas Law as it relates to one of the most American things ever: a lever-action rifle.

In our case, the lever-action rifle is chambered in .44 Remington Magnum and has a $20.00 in$ barrel.

This document refers to the  \href{https://drive.google.com/file/d/1JpZKagO_RQ2Az20Uw2yLxLpzfNrq7TlN/view?usp=sharing}{Module 12 Notes.pdf} document.


\section{Combustion of Gunpowder}
\begin{itemize}
\item one accepted equation for burning black powder is:\\
 \ce{2KNO3 (s) + S (s) + 3C (s) -> 3CO2 (g) + 2N (g) + K2S (g) +\Delta Q} 
\item see \cite{naval:gunpowder:combustion}
\item the chemistry of smokeless powders isn't quite so well documented, partially because there's not just one smokeless powder
\item so we're going to pretend this is the reaction fueling our lever-action rifle in .44 Remington Magnum
\footnote{Obviously loading a modern magnum cartridge with black powder isn't something you should try.}
\end{itemize}

\section{Pressure}

\begin{itemize}

\item pressure is force per unit area: $P = \frac{F}{A}$
\item pressure is measured in Pascals ($1 Pa = \frac{1 N}{1 m^{2}}$), in p.s.i ($1 psi = \frac{1 lb}{1 in^{2}}$), in atmospheres ($1 atm = 101 kPa$), etc.
\footnote{It honestly seems like everyone who has ever \emph{heard} of pressure has invented a new unit to measure it.}
\item the SAAMI spec has a .44 Remington Magnum running at a maximum pressure of  $250 MPa$ ($36000 psi$) 
\item we're going to assume our rifle is running at that maximum pressure
\footnote{In the real world, you're typically running loads significantly lower than the maximum load, hence you're running at lower than maximum pressure.}
\item notice how this compares to STP ($273 k$ and $101.325 kPa$, see Definition \ref{defn:stp}):
\begin{itemize}
\item $250 MPa$ is $250,000 kPa$
\item $250 MPa$ is $2467 atm$
\end{itemize}

\item that's a \emph{lot} of pressure
\end{itemize}

%\section{Ideal Gases}
%
%\begin{definition}[Ideal Gas]\label{def:ideal-gas}
%An ideal gas has the following properties:\\
%\begin{enumerate}
%\item The molecules or atoms that make up the gas are very small compared to the total volume available to the gas.
%\item The molecules of atoms that make up the gas must be so far apart from one another that there is no attraction or repulsion between them.
%\item The collisions that occur between the gas molecules or atoms must be elastic. When they collide, no energy can be lost in the collision.
%Also, no energy can be lost when they collide with the walls of the container.
%\end{enumerate}
%\end{definition}
%
%\begin{definition}[Standard Temperature and Pressure]\label{defn:stp}
%A temperature of 273k and a pressure of 1.00 atm.
%\end{definition}
%
%\begin{itemize}
%\item gases behave like ideal gases (see Definition \ref{def:ideal-gas}, p. \pageref{def:ideal-gas}) at high temperatures and low pressures
%\item we define the ``standard" temperature and pressure (STP) as 273k and 1.00 atm (see Definition \ref{defn:stp}, p, \pageref{defn:stp})
%\item so the gases in our rifle aren't really very ideal, but we're going to pretend they are pretty close (see ``lying to children")
%\footnote{In actual fact, the gun (and thus the gases in it) will heat up as we use it, so they'll be getting \emph{less} ideal.}
%\end{itemize}
%
%\pagebreak


\section{Ideal Gas Law}

\begin{law}[Ideal Gas Law]\label{law:ideal-gas}
$$P V = n R T$$ \\
$$ R = 0.0821 \frac{L \cdot atm}{mole \cdot K}$$
\end{law}

\begin{itemize}
%\item for an ideal gas (see Definition \ref{def:ideal-gas}, p. \pageref{def:ideal-gas}), we can related temperature, pressure, and volume by the Ideal Gas Law (Law \ref{law:ideal-gas}, p. \pageref{law:ideal-gas})
\item the Ideal Gas Law (Law \ref{law:ideal-gas}, p. \pageref{law:ideal-gas}) relates temperature, pressure, and volume
\footnote{The gases in a gun barrel aren't ideal gases -- not by a long shot -- and they'll become \emph{less} ideal as temperature increases as a result of shooting our gun. See ``lying to children."}
\item in that law:
\begin{itemize}
\item $n$ refers to the number of moles of our ideal gas
\item $R$ is a constant: $ R = 0.0821 \frac{L \cdot atm}{mole \cdot K}$
\item $P$ is pressure
\item $V$ is volume
\item $T$ is temperature
\end{itemize}
\end{itemize}

\section{Rifle Geometry}
\begin{itemize}
\item .44 Remington Magnum uses a $0.429 in$ bullet (typically $240 grains$)
\item if our rifle has a $20.00 in$ barrel, then we can calculate its volume:\\
\begin{equation} 
\begin{split}
 V &= \pi \cdot r^{2} \cdot h \\
    &= \pi (\frac{0.429 in}{2})^{2} (20.00 in) \\
    &= 2.8909 in^{3} \\
    &= 4.737 \times 10^{-5} m^{3} \\
    &= 4.437 \times 10^{-2} L \\
    &= 44.37 mL\\
\end{split}
 \end{equation}
\item so the lever-action rifle with a $20.00 in$ barrel has a barrel volume of $44 mL$
\item from the point of view of our chemical reaction, that's the \emph{maximum} volume of our gun barrel
\item the reaction occurs between the bullet and the cartridge base, so the starting volume is just the empty space in the cartridge
\item as the reaction occurs, the bullet moves down the barrel, and the volume increases
\footnote{Really the only difference between a rifle and a pipe bomb is that the bullet moves down the barrel of the rifle: if the bullet stops, then the rifle becomes a pipe bomb.
Many hunters have stories about this. Some have nicknames like ``Lefty."}
\item what's the theoretical maximum volume of the barrel for our cartridge?  it depends on $P V = n R T$, but it needs to be bigger than $44 mL$ for our rifle to work.
\footnote{By ``maximum volume," I mean it's possible to make a rifle barrel so long that the bullet just stops in there. 
This is exactly what a squib load is: a load where the theoretical maximum volume of the rifle barrel is smaller than the actual rifle barrel.}
\end{itemize}

\section{Shooting our Gun}
\begin{itemize}
\item when we squeeze the trigger on our lever gun, several steps happen very quickly:
\begin{enumerate}
\item the sear disengages, releasing the hammer
\item the hammer falls, striking the firing pin
\item the firing pin surges forward, striking the primer
\item the primer ignites, sending a spark/flame into the cartridge
\item the powder in the cartridge ignites, increasing pressure in the cartridge
\item the bullet moves forward into the barrel and hits the rifling
\end{enumerate}
\item things get interesting for us (as Chemists) when the powder ignites
\item it's easiest to contemplate this in terms of \emph{changes} that are all occurring:
\begin{enumerate}
\item $\Delta P$ is a change in pressure: initially $\Delta P \gg 0$ 
\item $\Delta V$ is a change in volume: $\Delta V > 0$ 
\item $\Delta n$ is a change in the number of moles of the gas(es) in the chamber: $\Delta n > 0$
\footnote{This isn't immediately obvious, but as the powder burns it consumes solids and produces gases. So $\Delta n > 0$ as long as our powder is burning.}
\item $\Delta T$ is a change in temperature: $\Delta T > 0$
\end{enumerate}
\item the most interesting change is $\Delta P$, because the manufacturer builds guns to handle a given pressure ($P$)
\item \emph{question} What happens if $P$ gets too large?
\item \emph{question} When would we expect $P$ to be the highest?
\item \emph{question} What happens if $\Delta n$ is too large? 
\item \emph{question} What happens if $\Delta T$ gets too large?
\item \emph{question} Why do longer gun barrels generally\footnote{Not always, see ``squib load" above.} mean faster bullets?
\end{itemize}

\section{Burn Rate}
\begin{itemize}
\item one important factor in choosing a powder is \emph{burn rate}
\item burn rate has a big effect on our $P$ because it's directly responsible for $\Delta n$
\item if we choose too fast a powder, then we'll end up with a very large $\Delta n$, which can cause too high a $\Delta P$
\item what do we mean by ``too high a $\Delta P$?" well, if the powder burns faster than the bullet can move, then we can get a ``pressure spike"
\item in other words, if our bullet moves too slowly, we can end up with a pipe bomb instead of a rifle
\footnote{Which is why every loading manual will have \emph{lower} powder charges for heavier bullets.}
\item this is a big consideration when loading for a magnum rifle
\item it's also a major difference between rifle and pistol powders
\footnote{Pistol and shotgun powders burn more quickly than rifle powders, and \emph{much} more quickly than magnum rifle powders.}
\item if you shoot a gun and you see a big ball of flame come out the muzzle there's a pretty good chance the powder charge is too large for that gun
\footnote{Of course with factory loads, this is way more likely than with hand loads.}
\end{itemize}

\section{Bullet Motion}
\begin{itemize}
\item the bullet moves as it's pushed by the expanding gases behind it
\item the force on the bullet depends on both pressure and the bullet's cross-sectional area
\item for our bullet, the cross-sectional area is:
\begin{equation} 
\begin{split}
 A &= \pi \cdot r^{2} \\
    &= \pi (\frac{0.429 in}{2})^{2} \\
    &= 0.1445 in^{2} \\
    &= 9.325 \times 10^{-5} m^{2} \\
\end{split}
 \end{equation}
\item we can calculate the force on the bullet from our definition of pressure:
\begin{equation} 
\begin{split}
P &= \frac{F}{A} \\
F &= P \times A \\
\end{split}
 \end{equation}
\item so our bullet moving under maximum pressure is being pushed by:
\begin{equation} 
\begin{split}
F &= P \times A\\
   &= (250 MPa) (9.325 \times 10^{-5} m^{2}) \\
   &= 23312.5 N\\
\end{split}
\end{equation}

\item or in English units:
\begin{equation} 
\begin{split}
F &= P \times A\\
   &= (36000 \frac{lb}{in^{2}})  (0.1445 in^{2}) \\
   &= 5202 lb \\
\end{split}
\end{equation}
 
\item that's a whole lot of push!
\footnote{There's a \emph{lot} of friction pushing back on the bullet from the barrel walls: $5000 lb$ seems like a lot of push (and it is), but there's a whole lot of resistance to overcome. }
\item the bullet will accelerate down the barrel until the pressure drops, usually because it leaves the gun barrel
\footnote{Again, see ``squib load" above.}
\item as the bullet moves, three very important quantities are changing:
\begin{enumerate}
\item $\Delta V > 0$  as the the space \emph{behind} the bullet increases
\item $\Delta T > 0$ as the powder burns (a highly exothermic reaction)
\item $\Delta n > 0$ as the chemical reaction produces more gas
\end{enumerate}

\item all of those changes will tend to drive pressure either up or down:
\begin{enumerate}
\item $\Delta P < 0$  because $\Delta V > 0$
\item $\Delta P > 0$  because $\Delta n > 0$
\item $\Delta P > 0$ because $\Delta T > 0$
\end{enumerate}

\item so pressure will be pushed \emph{both} up and down: whether it actually increases or decreases depends on the relative sizes of the other measurements
\footnote{With a low enough burn rate, for example, the pressure will \emph{decrease} as the bullet moves.}

\item when the bullet exits the barrel, the pressure behind it exits the barrel with a loud ``POP!" that sounds remarkably like a gunshot
\footnote{Because it \emph{is} a gunshot, obviously. Also, if the bullet is moving faster than sound there's a ``sonic boom" as well.}
\end{itemize}





\nocite{wile-chem-2}
\bibliography{../Chemistry}
\bibliographystyle{apalike}

\end{document}  